%% kapa: exam=1 round=34 year=2023 period=1 form=A subject=civil_law subject_ko=민법 seq=29 number=29
%% source: Q-Net 감정평가사 기출문제 (cst00309) artlSeq=5211789 file=raw/1008/5211789/2023년도 제34회 감정평가사 1차시험 1교시.pdf
%% conversion: kordoc 4.24.51 -> tools/gen1.py (auto); review=proofread-sonnet; flags=none
\begin{question}{29}
부동산의 점유취득시효에 관한 설명으로 옳지 않은 것은? (다툼이 있으면 판례에 따름)
\begin{choices}
  \item 집합건물의 공용부분은 취득시효에 의한 소유권 취득의 대상이 될 수 없다.
  \item 시효완성을 이유로 한 소유권취득의 효력은 점유를 개시한 때로 소급하지 않으며 등기를 함으로써 장래를 향하여 발생한다.
  \item 점유자가 점유 개시 당시에 소유권 취득의 원인이 될 수 있는 법률행위가 없다는 사실을 알면서 타인 소유의 토지를 무단점유한 것이 증명된 경우, 그 토지 소유권의 시효취득은 인정되지 않는다.
  \item 시효완성자는 취득시효의 기산점과 관련하여 점유기간을 통틀어 등기명의인이 동일한 경우에는 임의의 시점을 기산점으로 할 수 있다.
  \item 시효이익의 포기는 특별한 사정이 없는 한, 시효취득자가 취득시효완성 당시의 진정한 소유자에 대하여 하여야 한다.
\end{choices}
\end{question}
\answer{2}
